\documentclass[11pt, letterpaper]{article}

% --- Packages ---
\usepackage[margin=1in]{geometry} % Sets margins to 1 inch
\usepackage{fancyhdr}             % For custom headers and footers
\usepackage{amsmath, amssymb}     % For math equations and symbols
\usepackage{amsthm}               % For theorems and definitions (before newtx: \openbox clash)
\usepackage{newtxtext,newtxmath}
\usepackage{bm}
\usepackage{hyperref}             % For clickable links (TOC, URLs)
\usepackage[capitalise]{cleveref} % must follow hyperref
\usepackage{lastpage}             % To reference the total page count
\usepackage{natbib}
\usepackage{booktabs}
\usepackage{algorithm}
\usepackage{algpseudocode}   % from the algorithmicx package

% --- Header and Footer Setup ---
\pagestyle{fancy}
\fancyhf{} % Clear default header and footer

% EDIT HERE: Your details
\newcommand{\MyName}{T. Lucas Makinen}
\newcommand{\MyInstitution}{University of Cambridge}
\newcommand{\MyCourse}{Active dim degeneracy detection}
\newcommand{\MyDate}{\today}

% Header configuration
\lhead{\small \MyName \\ \MyInstitution} % Left Header: Name & Inst
\rhead{\small \MyCourse \\ \MyDate}      % Right Header: Course & Date

% Footer configuration
\cfoot{\thepage \ of \pageref{LastPage}} % Center Footer: Page X of Y

% Ensure header separation from text
\setlength{\headheight}{24pt}

% --- Custom Environments (Optional) ---
\newtheorem{theorem}{Theorem}
\theoremstyle{definition}
\newtheorem{definition}{Definition}
\newtheorem{example}{Example}
\newcommand{\matrixss}[1]{\ensuremath{\bm{\mathsf{#1}}}}

% --- Document Start ---
\begin{document}

% Title Block
\begin{center}
    {\Large \textbf{\MyCourse}} \\ [0.5em]
    % {\large \MyName \ | \ \MyInstitution} \\ [0.2em]
    {\small \MyDate}
\end{center}
\hrule
\vspace{1em}

% Table of Contents (Optional - comment out if not needed)
% \tableofcontents
\vspace{1em}
% \hrule
% \vspace{2em}

% --- Main Content Starts Here ---
\section*{Main Idea}

\section*{Amortised posterior in learned coordinates}

\paragraph{Setup.}
Parameters $\theta\in\Theta\subseteq\mathbb{R}^d$ with prior $\pi$, simulator
$x\sim p(\cdot\mid\theta)$. Fix an \emph{active} index set
$A\subset\{1,\dots,d\}$ with $|A|=m$ and complement $B$, and write
$\theta=(\theta_A,\theta_B)$. Learn two maps,
\[
  \eta_\phi:\mathbb{R}^d\to\mathbb{R}^m
  \quad\text{(coordinates, data--independent)},
  \qquad
  \hat\eta_\psi:\mathcal{X}\to\mathbb{R}^m
  \quad\text{(estimator)} .
\]

\paragraph{Model.}
Let $T(\theta)=\bigl(\eta_\phi(\theta),\,\theta_B\bigr)$ and
$J_A(\theta)=\partial\eta_\phi/\partial\theta_A\in\mathbb{R}^{m\times m}$.
Assume $T$ is a diffeomorphism on $\operatorname{supp}\pi$, which by
block-triangularity holds iff $\det J_A\neq0$. Define
\begin{equation}
  q_{\phi\psi}(\theta\mid x)
  \;=\;
  \mathcal{N}\bigl(\eta_\phi(\theta);\,\hat\eta_\psi(x),\,I_m\bigr)\;
  \pi_B(\theta_B)\;
  \bigl|\det J_A(\theta)\bigr| ,
  \label{eq:model}
\end{equation}
Here $\pi_B$ is the marginal prior of $\theta_B$. Substitute
$u=\eta_\phi(\theta_A;\theta_B)$ at fixed $\theta_B$ to get
$\int q\,d\theta_A=1$. Therefore $q$ integrates to one, and
$q(\theta_B\mid x)=\pi_B(\theta_B)$. The model states that $x$ is marginally
uninformative about $\theta_B$. It \emph{holds those components at the prior}.
For $m=d$ the set $B$ is empty, and \eqref{eq:model} is the usual change of
variables.

\paragraph{Objective.}
\begin{equation}
  \mathcal{L}(\phi,\psi)
  = -\,\mathbb{E}_{\theta\sim\pi,\;x\sim p(\cdot|\theta)}\bigl[\log q_{\phi\psi}(\theta\mid x)\bigr]
  = \mathbb{E}\Bigl[\tfrac12\bigl\|\eta_\phi(\theta)-\hat\eta_\psi(x)\bigr\|^2
    - \log\bigl|\det J_A(\theta)\bigr|\Bigr]
    + \tfrac{m}{2}\log 2\pi + H(\pi_B),
  \label{eq:loss}
\end{equation}
The term $H(\pi_B)$ is a constant. The objective inverts no matrix and needs no
Fisher matrix.

\paragraph{Proposition 1 (what the objective optimises).}
$\mathcal{L}=\mathbb{E}_x\bigl[\mathrm{KL}\bigl(p(\theta\mid x)\,\|\,q(\theta\mid x)\bigr)\bigr]+H(\theta\mid x)$,
so $\mathcal{L}\ge H(\theta\mid x)=H(\theta)-I(\theta;x)$, with equality if and
only if $q=p(\theta\mid x)$. The objective rules out the collapse
$\eta_\phi\equiv\mathrm{const}$, which sends $-\log|\det J_A|\to+\infty$.

\paragraph{Proposition 2 (relation to the Fisher target).}
Take $m=d$ and let $\theta^\star$ satisfy $\eta_\phi(\theta^\star)=\hat\eta_\psi(x)$.
Then $-\log q=\tfrac12(\theta-\theta^\star)^\top J^\top J(\theta-\theta^\star)+O(\|\cdot\|^3)$,
so $q$ has precision $J^\top J$. Matching the posterior precision $F$ gives
\begin{equation}
  J^\top J = F
  \quad\Longleftrightarrow\quad
  Q \;:=\; J^{-\top} F\, J^{-1} = I ,
\end{equation}
the two-stage flattening condition. The stationary points coincide. The route
to them does not need $J^{-1}$, $F^{-1}$, or $\operatorname{rank}F=d$.

\paragraph{Proposition 3 (uninformative directions, and the plateau).}
Suppose an active coordinate $\eta_i$ gets no information from $x$. The optimal
$\hat\eta_i$ is then constant. Minimise
$\mathbb{E}_{r}\bigl[\tfrac12(v(r)-c)^2-\log|v'(r)|\bigr]$ over scalar maps $v$.
The minimiser is the transport $v=\Phi^{-1}\!\bigl(F_r(r)\bigr)$. There the
induced density equals $p_r$, and the contribution to $\mathcal{L}$ is $H(p_r)$.
This equals exactly what the matching pass-through term contributes. An
uninformative active coordinate therefore changes $\mathcal{L}^\star$ by zero.

\paragraph{Corollary (intrinsic dimension).}
Let $r$ be the least $m$ such that $p(x\mid\theta)=p\bigl(x\mid s(\theta)\bigr)$
for some $s:\mathbb{R}^d\to\mathbb{R}^m$, and let
$\mathcal{L}^\star(m)=\min_{|A|=m}\min_{\phi,\psi}\mathcal{L}$. Then
$\mathcal{L}^\star$ is non-increasing. If the posterior is Gaussian in the
optimal coordinates, then
\begin{equation}
  \mathcal{L}^\star(m) > \mathcal{L}^\star(r) = H(\theta\mid x)
  \;\;\text{for } m<r,
  \qquad
  \mathcal{L}^\star(m) = \mathcal{L}^\star(r) \;\;\text{for } m\ge r .
\end{equation}
The plateau onset estimates $r$ from a likelihood rather than an eigengap.

\paragraph{Information spectrum.}
At the optimum the posterior covariance in $\eta$ is $I_m$. The prior covariance
eigenbasis is therefore canonical, and needs no Procrustes or varimax step:
\begin{equation}
  \lambda_i \;=\; \operatorname{Var}_\pi(\eta_i)
  \;=\; \frac{\operatorname{Var}_{\text{prior}}}{\operatorname{Var}_{\text{post}}},
  \qquad
  \tfrac12\log\lambda_i \;=\; \text{information gain of axis } i \text{ (nats)} .
\end{equation}
Uninformative axes give $\lambda_i=1$ by Proposition 3.

\paragraph{Worked example.}
$\mu(\theta)=(\theta_1,\;\theta_2-\theta_1^2,\;\theta_3,\;\theta_4)$,
$n$ i.i.d.\ replicates with $\varepsilon\sim\mathcal{N}(0,\operatorname{diag}(s_1^2,\dots,s_4^2))$,
$\pi=\mathcal{U}[-a,a]^4$, $w=2a$, and $s_3=s_4=S\to\infty$. The exact solution is
\[
  \eta_1=\frac{\sqrt n}{s_1}\theta_1,\qquad
  \eta_2=\frac{\sqrt n}{s_2}\bigl(\theta_2-\theta_1^2\bigr),\qquad
  \det J_A=\frac{n}{s_1s_2},\qquad
  \lambda=\Bigl(\tfrac{n a^2}{3s_1^2},\;\tfrac{n}{s_2^2}\bigl(\tfrac{a^2}{3}+\tfrac{4a^4}{45}\bigr),\;1,\;1\Bigr).
\]
With $A=\{1,2\}$,
\begin{equation}
  \mathcal{L}^\star(2)=1-\log\frac{n}{s_1s_2}+\log 2\pi+2\log w,
  \qquad
  \mathcal{L}^\star(1)=\mathcal{L}^\star(2)+\log w-\tfrac12+\log\frac{\sqrt n}{s_2}-\tfrac12\log 2\pi .
\end{equation}
Set $n=50$, $s=(1,2,\infty,\infty)$ and $a=3$. Then
$\lambda=(150,\,127.5,\,1,\,1)$, and
$\mathcal{L}^\star=(4.838,\,3.203,\,3.203,\,3.203)$ nats per sample for
$m=1,2,3,4$. The gain from $m=1$ to $m=2$ is $1.636$, and every later gain is
zero, so $r=2$.

\paragraph{Remark (sub-rank models are conditional).}
The value $\mathcal{L}^\star(1)=4.838$ above is the optimum \emph{restricted} to
maps $\eta_\phi=\eta_\phi(\theta_A)$. Equation \eqref{eq:model} lets $\eta_\phi$
read all of $\theta$. Its first factor is then $q(\theta_A\mid\theta_B,x)$, not
$q(\theta_A\mid x)$. For $m<r$ the model can therefore exploit dependence
between $\theta_A$ and $\theta_B$ under the likelihood. Take the worked example
with $A=\{1\}$. Conditioning on $\theta_2$ makes $\bar x_2$ informative about
$\theta_1^2$, which raises the Fisher information for $\theta_1$ from $n/s_1^2$
to $n/s_1^2+4n\theta_1^2/s_2^2$. Hence
\begin{equation}
  \mathcal{L}^\star(1)\;\ge\;4.838-\tfrac12\,\mathbb{E}_\pi
  \log\!\Bigl(1+\tfrac{4s_1^2}{s_2^2}\theta_1^2\Bigr)
  \;=\;4.838-0.568\;=\;4.271 .
\end{equation}

Only part of this bound is reachable. The estimator $\hat\eta_\psi$ enters as a
single location shift. It therefore cannot reproduce how the true conditional
depends on both of its arguments. The $d=20$ experiment below reaches $0.332$
of the $0.568$ bound. Two consequences follow. First, the plateau \emph{location} does not
move, because Proposition 3 concerns $m\ge r$, and
$\mathcal{L}^\star(m)>\mathcal{L}^\star(r)$ still holds strictly for $m<r$.
Second, the measured gain $\mathcal{L}^\star(m)-\mathcal{L}^\star(m+1)$
\emph{underestimates} the information in the promoted axis. Read a gain as
per-axis information only when $\eta_\phi$ reads $\theta_A$ alone.

\section*{Estimation}

\paragraph{Choosing the active set: why $A$ is a coordinate subset.}
Equation \eqref{eq:model} needs $\pi_B$ to be a tractable marginal that is
independent of the active block. Write a general orthonormal split
$\theta=R_1u+R_2v$ with $R=[R_1\,R_2]\in O(d)$. That independence holds for
every $R$ if and only if $\pi$ is rotationally invariant. A Gaussian prior
$\pi=\mathcal{N}(0,I_d)$ qualifies. A learned frame
$V\in\mathbb{R}^{d\times m}$ of full column rank then replaces $A$, and gives
\begin{equation}
  -\log q = \tfrac12\bigl\|\eta_\phi(\theta)-\hat\eta_\psi(x)\bigr\|^2
  + \tfrac12\Bigl(\|\theta\|^2-\theta^\top V(V^\top V)^{-1}V^\top\theta\Bigr)
  - \log\bigl|\det (JV)\bigr| + \tfrac12\log\det(V^\top V)
  + \tfrac d2\log2\pi ,
  \label{eq:gaussV}
\end{equation}
where $J=\partial\eta_\phi/\partial\theta$. Gradient descent learns the
subspace, and the cost of discovery does not grow with $d$.

Two practical warnings apply. First, leave $V$ unconstrained and use
\eqref{eq:gaussV} as written. A polar retraction that maps $V$ onto the Stiefel
manifold through $V(V^\top V)^{-1/2}$ makes $V^\top V=I_m$ exactly. Its
eigendecomposition then has repeated eigenvalues and a singular derivative. The
gradients are therefore not finite at the start of training.

Second, the construction fails outright for $\pi=\mathcal{U}[-a,a]^d$.
Here $v=R_2^\top\theta$ projects a
hypercube onto a subspace, which gives a marginal of Irwin--Hall type. That
marginal has no closed form, and it depends on $u$. The term $H(\pi_B)$ in
\eqref{eq:loss} is therefore neither constant nor computable, and a free $V$
would tilt to exploit $H(\pi_A)+H(\pi_B)>H(\pi)$. Under a box prior the set $A$
must be a coordinate subset, and \eqref{eq:loss} reduces to
$\mathcal{L}=\mathbb{E}[\tfrac12\|\cdot\|^2-\log|\det J_A|]+\tfrac m2\log2\pi+(d-m)\log w$.

\paragraph{Screening.}
A search over coordinate subsets costs $\binom dm$ fits. A screen avoids that
cost. The informative coordinates are exactly those on which the likelihood
depends. For
a model with sufficient statistic $t(x)$, those are the coordinates that enter
$\mathbb{E}[t\mid\theta]$. Fit one regression
$g_\chi:\mathbb{R}^d\to\mathbb{R}^{\dim t}$ with a group penalty on the columns
of its first weight matrix $W^{(1)}\in\mathbb{R}^{d\times h}$,
\begin{equation}
  \min_\chi\;\mathbb{E}\bigl\|t(x)-g_\chi(\theta)\bigr\|^2
  \;+\;\lambda\sum_{j=1}^{d}\bigl\|W^{(1)}_{j,:}\bigr\|_2 ,
  \label{eq:screen}
\end{equation}
then rank the coordinates by $\|W^{(1)}_{j,:}\|_2$.

The screen uses no prior,
and costs one fit rather than $\binom dm$. It returns a \emph{ranking}, so the
sets $A_m=\{$top $m\}$ nest, and the plateau of the Corollary still selects
$m$. The screen finds coordinates that enter nonlinearly or through
interactions. It fails only for a coordinate that changes higher moments of $t$
but leaves $\mathbb{E}[t\mid\theta]$ constant. Append a higher-moment summary
to $t$ in that case.

\paragraph{Calibration.}
The model makes two claims that held-out pairs can falsify without an oracle.
The first claim is
$z:=\eta_\phi(\theta)-\hat\eta_\psi(x)\sim\mathcal{N}(0,I_m)$, exactly. Check
the moments of $z$, its coverage at $|z|<1,1.96,2.58$, and the uniformity of
$\Phi(z_i)$. A standard deviation above one means the posterior is narrower
than the data supports. The second claim is that $\theta_B$ follows $\pi_B$ and
does not depend on $x$. Check the $\pi_B$-PIT for uniformity.

The sharper test regresses each $\theta_j$, $j\in B$, on a sufficient summary of
$x$. Compare the resulting $R^2$ with its null value $k/N$, for $k$ regressors
and $N$ test points. A leakage $R^2$ well above $k/N$ says that $m$ is too
small. It says so per coordinate, and it does not use the plateau. Agreement
between the two diagnostics is stronger evidence than either one alone.

\paragraph{Uncertainty.}
At the optimum the posterior covariance in $\eta$ is $I_m$. An ensemble spread
therefore arrives in units of the statistical error, and needs no normalisation
convention. Train $K$ members on resampled pairs $(\theta,x)$. Use the
out-of-bag pairs of each member for early stopping, which costs no extra
simulations. Fix the $O(m)$ gauge of every member to a common frame, and report
the per-point spread $\sigma_{\mathrm{epi},i}$. This spread is the
$y_{\mathrm{std}}$ for a later symbolic regression. The factor
$1+\sigma_{\mathrm{epi},i}^2$ inflates the effective posterior variance to
account for map uncertainty.

Check the spread with the leave-one-out statistic
\[
  z_k=\bigl(\eta_k-\bar\eta_{-k}\bigr)\big/\bigl(s_{-k}\sqrt{K/(K-1)}\bigr) .
\]
For exchangeable
Gaussian members this statistic follows $t_{K-2}$. The reference standard
deviation is therefore $\sqrt{(K-2)/(K-4)}$, not $1$. Run the screen
\eqref{eq:screen} inside each replicate as well. This gives the bootstrap
frequency at which the screen recovers $A$, a discrete analogue of subspace
uncertainty.

\paragraph{Scoring symbolic candidates.}
To score a candidate coordinate map $g:\mathbb{R}^d\to\mathbb{R}^{m}$, freeze
$\eta_\phi(\theta)=s\odot\bigl(\Lambda\,g(\theta)\bigr)$. Keep the matrix
$\Lambda\in\mathbb{R}^{m\times m}$ and the diagonal $s$ learnable. The score
then judges $g$ up to the $O(m)$ gauge and the scale. Refit $\psi$, then
evaluate
$\hat{\mathcal{L}}$ on held-out pairs. The two-part description length is
\begin{equation}
  \mathrm{DL}(g) \;=\; \mathrm{DL}_{\mathrm{model}}(g)
  \;+\;\sum_{i=1}^{N}\bigl(-\log q(\theta^{(i)}\mid x^{(i)};g)\bigr),
  \label{eq:mdl}
\end{equation}
in nats throughout. The \emph{paired} per-sample differences separate the
candidates, and their standard error is far smaller than that of either term
alone.

Two properties matter here. Equation \eqref{eq:model} is a normalised density on
$\mathbb{R}^d$ for every $m$, so candidates of different output dimension
compare directly. One scale therefore selects both the rank and the functional
form. The data term of \eqref{eq:mdl} grows with $N$, but
$\mathrm{DL}_{\mathrm{model}}$ does not. Complexity discriminates only when
per-sample fit differences drop below $\mathrm{DL}_{\mathrm{model}}/N$. Always
quote $N$ with a description length.

\paragraph{An over-specified $m$ is safe.}
A sweep from $m=1$ is not necessary. Proposition 3 fixes the value that an
uninformative active coordinate takes. It settles at the transport of its own
prior marginal, so $\lambda_i=1$ exactly. The spectrum therefore has a known
null value. One fit at a generous $m$ gives $r$.

Table~\ref{tab:overspec} tests this claim at $m=2,3,5,8$ against a true rank of
$2$. Two eigenvalues stay far above $1$ in every case. The others land within
$0.06$ of $1$. The top two canonical axes recover the true coordinates to
$R^2=0.999$. The reverse regression returns the same figure, so the learned
plane and the true pair are mutually invertible.

The discarded axes hold almost nothing about the truth, at $R^2\le0.03$. Against
the spurious screened coordinates they reach $R^2$ between $0.96$ and $0.99$.
The model spends a spare dimension on the transport of a nuisance parameter.
That is the optimum, not a failure.

\paragraph{The price of a spare dimension.}
Each spare dimension costs about $0.085$ nats under the box prior. The reason is
specific to the box. The best uninformative coordinate is
$\eta=\Phi^{-1}\bigl((\theta+a)/w\bigr)$. This function grows without bound at
the faces of the box. A finite network with a linear skip cannot fit those
tails.

Under $\pi=\mathcal{N}(0,I_d)$ the same transport is the identity map and costs
nothing. The Gaussian run paid $0.0038$ nats for one spare dimension, against
$0.0903$ here. The penalty brings one benefit and one risk. It tilts the plateau
upward, so $\arg\min_m\hat{\mathcal{L}}$ gives $r$ with no tolerance parameter.
But a direction that carries less than $0.085$ nats hides inside the penalty.
The spectrum still shows that direction as $\lambda_i>1$.

Read the rank from the spectrum. Use $\hat{\mathcal{L}}$ only to confirm it.

\paragraph{Protocol when $r$ is unknown.}
\begin{enumerate}\itemsep0pt
  \item Rank all $d$ coordinates with the screen \eqref{eq:screen}.
  \item Fit one model at a generous $m$, about $1.5$ times the expected rank.
  \item Count the eigenvalues above the null band. Set $\hat r$ to that count.
  \item Refit at $\hat r$. Check that $\hat{\mathcal{L}}$ decreases by about
        $0.085\,(m-\hat r)$.
  \item Check the $z$ calibration and the complement leakage at $\hat r$.
  \item Bootstrap $K$ members. Read the error bars on $\lambda_i$ and the
        recovery frequency of $A$.
  \item Score the symbolic candidates by \eqref{eq:mdl} with paired tests.
\end{enumerate}
Steps 1 and 2 replace the sweep over $m$. Two fits and one bootstrap are enough.
A generous $m$ also protects the screen. The active set is the top $m$ of the
ranking. A wide net keeps an informative parameter that the screen ranks a
little too low.

Every step reuses the same simulated pairs. Only step 2 depends on $m$. The
geometry terms of \eqref{eq:loss} involve $\theta$ alone, so fresh prior draws
evaluate them at no simulation cost. Draw those samples from $\pi$. Any other
law trains the volume term and the residual term on inconsistent supports.

\paragraph{The null band on $\lambda$.}
Sampling alone lifts the largest spurious eigenvalue to $(1+\sqrt{m'/N})^2$, for
$m'$ spare dimensions and $N$ test points. This gives $1.035$ at $m'=6$ and
$N=20{,}000$. The measured value was $1.055$. Transport error sets the band, not
sampling noise. Take the band from the bootstrap spread of the near-unit
eigenvalues.

\paragraph{Graded spectra.}
Report the information profile $\tfrac12\log\lambda_i$ in nats. Do not report a
single rank. The demonstration below has an extreme gap, from $194$ down to
$1.05$. A real problem with $7$ to $10$ informative directions gives a graded
spectrum. An example is $500,200,50,12,3,1.4,1.1,1.02$. The rank is then not a
well posed quantity.

In nats the question becomes how much information to discard. A reader can check
that choice. A threshold on an eigengap gives no such handle. Two close
eigenvalues also leave their own plane unidentified. Run the symbolic regression
on such a pair together.

\paragraph{Irrelevant parameters against irrelevant directions.}
The box prior needs a coordinate subset $B$ that the data leaves at the prior.
It can express an irrelevant parameter. It cannot express an irrelevant
direction. Suppose the data constrains every parameter, but only $m$ of their
combinations carry information. No subset $B$ then exists, and the screen keeps
all $d$ coordinates.

Use the learned frame $V$ of \eqref{eq:gaussV} for that case. Transform the
prior into a standard normal first. A spare dimension is then free, because the
transport is the identity map. A subspace complement also drops the
axis-alignment restriction on $A$. A problem with $d=30$ and a suspected rank
near $8$ usually sits in this class.

\section*{Numerical demonstration}

Take $d=20$ with
$\mu(\theta)=\bigl(\theta_{i_0},\,\theta_{i_1}-\theta_{i_0}^2\bigr)$ for one
hidden pair $(i_0,i_1)$. Use $n=8$ replicates, $s=(0.25,0.5)$,
$\pi=\mathcal{U}[-3,3]^{20}$, and $32{,}000$ training and $20{,}000$ test pairs.
Then $18$ of the $20$ parameters are pure nuisance, and the two that matter are
degenerate. The screen \eqref{eq:screen} gives group norms $0.344$ and $0.067$
for $\theta_{i_0}$ and $\theta_{i_1}$, against $\le0.004$ for the other $18$.
It isolates the pair at once.

\begin{table}[h]
\centering
\begin{tabular}{cccl}
\toprule
$m$ & $\hat{\mathcal{L}}(m)$ & restricted $\mathcal{L}^\star(m)$ & \\
\midrule
1 & 32.704 & 33.036 & below the restricted value (see Remark) \\
2 & 30.944 & 30.931 & $\hat r=2$; correct coordinate set \\
3 & 31.035 & 30.931 & plateau \\
\bottomrule
\end{tabular}
\caption{Rank sweep, nats per sample on held-out pairs.}
\end{table}

\begin{table}[h]
\centering
\begin{tabular}{cccll}
\toprule
$m$ & $\hat{\mathcal{L}}$ & excess & spectrum $\lambda_i$ & $R^2$ for the true pair \\
\midrule
2 & 30.944 & $+0.014$ & $507,\, 265$ & $0.9998,\, 0.9997$ \\
3 & 31.035 & $+0.104$ & $593,\, 199,\, 1.02$ & $0.9997,\, 0.9994$ \\
5 & 31.174 & $+0.244$ & $682,\, 149,\, 1.05,\, 1.02,\, 1.00$ & $0.9998,\, 0.9992$ \\
8 & 31.461 & $+0.531$ & $571,\, 194$, then six in $[0.98,1.06]$ & $0.9998,\, 0.9991$ \\
\bottomrule
\end{tabular}
\caption{An over-specified $m$ against a true rank of $2$. The excess column
gives the gap to the floor $\mathcal{L}^\star(2)=30.931$, and it grows by about
$0.085$ nats per spare dimension. The last column regresses each true coordinate
on a cubic in the top two canonical axes.}
\label{tab:overspec}
\end{table}

Now check calibration at $\hat r=2$. The residual $z$ has means $-0.019$ and
$-0.007$, and standard deviations $1.040$ and $1.017$. Coverage at $|z|<1.96$ is
$0.941$ and $0.946$, against a nominal $0.950$. The model is therefore
overconfident by $2$ to $4$ percent. The complement PIT has standard deviation
$0.2890$ against the uniform value $0.2887$.

Complement leakage is $\overline{R^2}=0.00020$ against the null value
$4/20{,}000=0.00020$. The test therefore confirms, as far as it can resolve,
that the data says nothing about the other $18$ directions. All $8$ bootstrap
replicates re-screen to the true pair on their own. The leave-one-out $z$
reaches $0.918$ and $0.971$ of the $t_6$ reference. The ensemble spread is
therefore honest to a few percent.

\begin{table}[h]
\centering
\begin{tabular}{lccrr}
\toprule
candidate & $m$ & $\mathrm{DL}_{\mathrm{model}}$ & $\hat{\mathcal{L}}$ & vs.\ neural \\
\midrule
$\theta_{i_0},\;\theta_{i_1}-\theta_{i_0}^2$      & 2 & 9  & 30.891 & $-0.053\pm0.003$ \\
$\theta_{i_0},\;\theta_{i_1}-\theta_{i_0}^{1.9}$  & 2 & 14 & 30.905 & $-0.040\pm0.003$ \\
$\theta_{i_0},\;\theta_{i_1}$                     & 2 & 4  & 31.516 & $+0.572\pm0.007$ \\
$\theta_{i_0},\;\theta_{i_1}-\theta_{i_0}^2-0.3\theta_j$ & 2 & 15 & 32.041 & $+1.097\pm0.006$ \\
$\theta_{i_0}$                                    & 1 & 2  & 33.012 & $+2.067\pm0.006$ \\
\midrule
neural $\eta_\phi$ (reference)                    & 2 & --- & 30.944 & --- \\
\bottomrule
\end{tabular}
\caption{Candidate scoring, nats per sample, $N=20{,}000$ held-out pairs.
Rows are ordered by \eqref{eq:mdl}.}
\end{table}

Three features of the table carry the argument. First, the true expression
\emph{beats} the learned map by $0.053\pm0.003$ nats, or $19$ standard errors.
The symbolic answer is a strictly better model of the posterior than the network
that proposed it. That claim is stronger than a flatness residual, and it also
validates the network.

Second, the exponents $2.0$ and $1.9$ differ by only $0.013\pm0.001$ nats per
sample. The direct paired test resolves this gap at $9$ standard errors. A
comparison of each candidate against the network does not. Pairwise comparison
is what separates near-ties.

Third, the spurious nuisance term costs more ($1.097$) than dropping the
quadratic altogether ($0.572$). An irrelevant coordinate injects prior variance
that the data cannot predict. The penalty is
$\tfrac12\log\bigl(1+0.3^2\operatorname{Var}_\pi(\theta_j)n/s_2^2\bigr)=1.13$,
which matches the measurement. The objective therefore rejects spurious
variables with no sparsity prior at the scoring stage. It rejects them more
sharply under a broad prior than a narrow one.

One number needs comment: $30.891<30.931$. The closed-form $\mathcal{L}^\star$
uses the differential entropy of an untruncated Gaussian for the active block.
The estimator $\hat\eta_\psi$ instead learns the posterior-mean shrinkage near
the faces of the box. Under a box prior, expect a few hundredths of a nat below
the stated floor. Under a Gaussian prior any such dip indicates overfitting.


\bibliographystyle{mnras}
\bibliography{main}

\end{document}